\section{系统模型}

\subsection{参与实体}

系统包含五类实体。

\textbf{数据提供方}负责注册数据对象、生成分片、计算哈希承诺和设置审计策略。\textbf{链下存储节点}保存数据分片、校验分片或计算证据，并响应审计挑战。\textbf{审计节点}发起抽样挑战，检查分片哈希或证明结果，并提交审计摘要。\textbf{审计委员会}由高可信节点组成，负责对审计摘要、恢复状态和信誉变化进行联盟链确认。\textbf{联盟链}保存数据对象状态、Merkle 根、哈希链锚点、恢复事件和节点信誉。

\subsection{链上链下分工}

TrustFlowAudit 采用链下计算加链上存证的分层结构。链下层负责数据切片、纠删码编码、分片保存、抽样校验、异常恢复和完整证据保存；链上层负责保存不可篡改的摘要证据，包括数据对象标识、编码参数、分片 Merkle 根、审计批次摘要、恢复状态和信誉变化。

图~\ref{fig:architecture} 给出系统架构。该架构强调链上不直接承担高频证明计算，而是作为可信时间线和责任账本使用。

\begin{figure}[t]
\centering
\begin{tikzpicture}[
  node distance=0.55cm,
  box/.style={draw, rounded corners, align=center, minimum width=2.0cm, minimum height=0.65cm, font=\footnotesize},
  wide/.style={draw, rounded corners, align=center, minimum width=6.8cm, minimum height=0.75cm, font=\footnotesize},
  arrow/.style={-{Latex[length=2mm]}, thick}
]
\node[box] (owner) {数据提供方};
\node[box, right=0.5cm of owner] (storage) {链下存储节点};
\node[box, right=0.5cm of storage] (auditor) {审计节点};
\node[wide, below=0.65cm of storage] (offchain) {链下层：分片、哈希校验、纠删码恢复、证据保存};
\node[wide, below=0.65cm of offchain] (onchain) {链上层：Merkle 根、审计摘要、恢复状态、信誉更新};
\node[box, below=0.65cm of onchain] (committee) {可信审计委员会};
\draw[arrow] (owner) -- (storage);
\draw[arrow] (storage) -- (auditor);
\draw[arrow] (owner) -- (offchain);
\draw[arrow] (storage) -- (offchain);
\draw[arrow] (auditor) -- (offchain);
\draw[arrow] (offchain) -- (onchain);
\draw[arrow] (committee) -- (onchain);
\end{tikzpicture}
\caption{TrustFlowAudit 链下计算与链上存证架构}
\label{fig:architecture}
\end{figure}

\subsection{数据对象与承诺}

设数据对象 $D$ 被切分为 $k$ 个数据分片，并生成 $m$ 个校验分片。每个分片 $s_i$ 的哈希记为 $h_i=H(s_i)$。系统根据所有分片哈希构造 Merkle 根 $R_D$，并将 $\langle id_D,k,m,R_D,policy\rangle$ 写入链上状态，其中 $policy$ 表示抽样比例、审计周期、风险等级和恢复策略。

链下证据按照哈希链组织。对于审计事件 $e_t$，系统计算：
\begin{equation}
  c_t = H(c_{t-1} \parallel e_t),
\end{equation}
其中 $c_t$ 是第 $t$ 个事件后的哈希链状态。链上只需要周期性保存 $c_t$ 或其批量 Merkle 根，即可证明链下证据序列未被无痕修改。
